\section{Supervision for the L11 extension}
The optional L11 selector additionally uses independent text/image/category
utility supervision. Its 22 fit, three calibration and four development rooms
provide 629/207/175 text expressions, 237/28/25 image queries and 116 category
episodes. Only 88 fit-room category episodes enter training; calibration and
development use full maps. Query-image views and neighbors are excluded from
their candidate observations; full-scan geometry remains available. One margin
(0.15052) protects 15 calibration room/output cells. Five room folds fit a
64-tree structural prior plus a 64-tree query residual. All 19 ScanNet/UQIS jobs
are predicted before scoring. These are extension settings, not prerequisites
for the v10 main results. Development and calibration informed model design.

\section{Measured optional replacement controls}

All rows in Table~\ref{tab:ablations} retain the same teachers, complete
candidates, supervision rooms and output definitions. The query-blind arm returns only the fitted structural prior (64 trees,
versus 128 in the full model); it therefore also changes capacity. Permutation
instead uses the fitted full model and
shuffles all teacher scores jointly within each view. It preserves their
multisets and correlations. Removing the null constraint retains the same
64+64-tree capacity and total fit weight, replacing permuted training examples
with duplicates of aligned examples. Each trained control uses the same
calibration procedure. The zero-margin control uses the full fitted model and
changes only the acceptance margin. These are response-level ablations.

\begin{table}[ht]
\centering\small\resizebox{\linewidth}{!}{\begin{tabular}{llrrrrrrr}
\toprule
Policy & Repair / harm & S19 & S15 & S10 & T-IoU & I-IoU & T-AP & I-AP \\
\midrule
Incumbent only & 0 / 0 & 61.01 & 62.53 & 72.49 & 16.39 & 39.90 & 26.60 & 59.31 \\
Full shared replacement & 1 / 0 & 61.01 & 62.53 & 72.49 & 16.39 & 40.99 & 26.60 & 60.58 \\
Without query coordinates & 1 / 1 & 61.01 & 62.53 & 72.49 & 16.39 & 40.26 & 26.60 & 59.50 \\
Permuted query correspondence & 0 / 0 & 61.01 & 62.53 & 72.49 & 16.39 & 39.90 & 26.60 & 59.31 \\
Without null correction constraint & 1 / 0 & 61.01 & 62.53 & 72.49 & 16.39 & 40.99 & 26.60 & 60.58 \\
Without positive margin & 56 / 33 & 60.54 & 62.05 & 71.62 & 22.17 & 41.17 & 35.18 & 60.53 \\
\bottomrule
\end{tabular}
}
\caption{Matched controls on all ten ScanNet rooms and 268 UQIS queries.
S19/15/10 are mesh mIoU. Repair/harm counts concern actual UQIS set IoU;
unchanged actual points are included in the denominator. This is one coupled
policy comparison, not a collection of per-benchmark best runs.}
\label{tab:ablations}
\end{table}

The image gain over query-blind and permuted controls is respectively 0.725 and
1.088 points; descriptive nine-room bootstrap intervals are [0,2.174] and
[0,3.264]. They do not establish a robust query-specific advantage. The full
and no-null arms make identical choices in all 144 text/image/menu decisions,
so the null constraint has no demonstrated accuracy contribution. Removing the
positive margin improves text and image IoU, but reduces ScanNet mIoU, including
a 0.869-point loss on the ten-class menu. Retention prevents some damage while
also rejecting useful instance replacements. Neither the full policy nor the
zero-margin variant is promoted to replace the reference system.

\section{Candidate capacity and failure attribution}
For a fixed candidate pool, let $G_0$ be gain under zero-margin selection and
$G_\tau$ gain under calibrated retention. The exact diagnostic is
\[
G_{\rm pool}=R_{\rm rank}+(G_0-G_\tau)+G_\tau.
\]
The middle term is signed: it is negative when retention prevents damage.
This identity prevents attributing all unused candidate capacity to the margin.
For text, 15.395 points of available IoU gain decompose into 9.615 ranking regret,
5.780 retention cost and zero actual gain. For images, 9.492 available points
decompose into 8.224 ranking regret, 0.180 retention cost and 1.088 actual gain.
Thus image ranking is the dominant measured bottleneck; lowering the margin
cannot recover most of the unused capacity. Category mean-menu capacity is
only 0.913 points, and its negative retention cost indicates prevented damage.
These label-oracle diagnostics are computed after prediction and are never
automatic method results.


\begin{center}
\resizebox{\linewidth}{!}{\begin{tabular}{llrrrrrrrr}
\toprule
Method & Origin & T-IoU & T-AP & I-IoU & I-AP & P2-IoU & P2-AP & P3-IoU & P3-AP \\
\midrule
Ours (regional readout) & O & 12.11 & 23.69 & 26.55 & 37.75 & 42.09 & 52.74 & 41.52 & 52.36 \\
Ours (v10) & O & 16.39 & 26.60 & 39.90 & 59.31 & 49.74 & 75.79 & 49.05 & 75.33 \\
Ours + replacement (L11) & M & 16.39 & 26.60 & 40.99 & 60.58 & 49.74 & 75.79 & 49.05 & 75.33 \\
\bottomrule
\end{tabular}
}
\end{center}
The complete UQIS comparison retains the established method and the optional
selector as separate rows. The selector does not explain the established
method's historical component gains.
